% Independent contribution. AI-produced by GPT-6 Astra; no priority claimed.
\subsection{The literal descent map and its actual kernel}

Write $G=\widehat{GT}_0$ and retain R1's faithful realization as
the outer centralizer of the geometric $\mathfrak S_3$ acting
on the actual pure subgroup $F\triangleleft B$, where
\[
 B=\langle s,r\mid s^2=r^3=1\rangle^{\wedge},\quad
 \tau=sr,\quad x=\tau^{-2},\quad y=sxs,\quad z=(xy)^{-1}.
 \tag{B.C1}
\]
The completed free group $F=\widehat F(x,y)$ has index six.
The map in the literal question is
$\varphi_P:F\twoheadrightarrow B$, $x\mapsto s$, $y\mapsto\tau$.
Its kernel is the closed normal closure of $x^2,z^3$.
The normalized source action takes each of these peripheral
classes to a conjugate of a unit power; the inverse action
does the same.  Thus it preserves this kernel, and descent
is an actual continuous automorphism of $B$, up to inner
automorphism.  Restriction to the characteristic pure subgroup,
followed by R1's normalization, defines a continuous
endomorphism $D_P:G\to G$.  The same construction applies to
the quotient $\varphi_L$ specified below.

\begin{result}{B.1.1}{Proved}{The exact conversion of the two descent markings}
The LM2 generators and physical quotient satisfy
\[
 x_L=(rs)^2=y^{-1},\qquad y_L=(sr)^2=x^{-1},\qquad
 \varphi_L(x_L)=s,\quad\varphi_L(y_L)=\tau.
\]
Equivalently, on the current physical subgroup,
\[
 \varphi_L(x)=\tau^{-1},\quad\varphi_L(y)=s,
 \qquad\varphi_L=\varphi_P\kappa,\quad
 \kappa(x)=y^{-1},\quad\kappa(y)=x^{-1}.
 \tag{B.C2}
\]
For $c=(-1,1)\in G$, their actual descent endomorphisms obey
\[
 D_L(g)=D_P(cgc)=cD_P(g)c,\qquad D_P(c)=c.                 \tag{B.C3}
\]
In particular the two precisely defined fixed subgroups are
\[
 C_P=\operatorname{Fix}(D_P),\qquad
 C_L=\operatorname{Fix}(D_L)=\{g:D_P(g)=cgc\}.
 \tag{B.C4}
\]
The $C_\varphi$ in the literal current question is $C_P$.
\end{result}
\begin{proof}
The word substitutions in (B.C2) prove all four values on
free generators.  In $\operatorname{Aut}(F)$,
$\kappa=\sigma c=c\sigma$, where $\sigma(x)=y$ and
$\sigma(y)=x$.  For any quotient and any automorphism
preserving its kernel, changing the quotient from $\varphi_P$
to $\varphi_P\kappa$ changes the descended class from that
of $g$ to that of $\kappa g\kappa^{-1}$.  Every $g\in G$
commutes with the geometric interchange $\sigma$ in the
outer group.  Hence this last class is $cgc$, giving the first
equality of (B.C3).

The descent of $c(x)=x^{-1},c(y)=y^{-1}$ through
$\varphi_P$ is represented on $B$ by
\[
 s\mapsto s,\qquad \tau\mapsto\tau^{-1},\qquad
 r\mapsto sr^{-1}s.
 \tag{B.C5}
\]
On $B/F=\mathfrak S_3$ it is the identity; its restriction
sends the named $x,y$ to their inverses.  Thus $D_P(c)=c$.
The homomorphism law for descent now gives the other equality
of (B.C3).  This proof retains both quotient maps and the
automorphism relating them.
\end{proof}

\begin{result}{B.1.2}{Proved}{Arithmetic descent and a degree-sixteen counterexample to the literal inclusion}
On the arithmetic image,
\[
 D_L(\gamma)=\gamma,\qquad
 D_P(\gamma)=c\gamma c\quad(\gamma\in G_{\mathbb Q}),
 \qquad C_P\cap G_{\mathbb Q}=\{1,c\}.
 \tag{B.C6}
\]
There is an actual $\gamma\in G_{\mathbb Q^{\rm ab}}$
for which $D_P(\gamma)\ne\gamma$, already distinguished by
the Galois extension $\mathbb Q(\zeta_8,2^{1/8})/\mathbb Q$.
Consequently the literal inclusion
$\widehat{GT}\subseteq C_\varphi$ and both proposed
equalities $C_\varphi=\widehat{GT}$ and
$C_\varphi=\ker\Delta$ are refuted.
\end{result}
\begin{proof}
Here is the actual arithmetic comparison, including loop
orientation.  Let $\mathcal X/\mathbb Q$ have coarse curve
$\mathbb P^1\setminus\{0\}$, with inertia orders two at
$1$ and three at $\infty$, in coordinate $T=1728/j$.
The two maps from $U=\mathbb P^1\setminus\{0,1,\infty\}$
are
\[
 q(t)=\frac{27t^2(1-t)^2}{4(1-t+t^2)^3},\qquad
 j:U\hookrightarrow\mathcal X,\quad j(t)=t.
 \tag{B.C7}
\]
They are the Legendre--Jacobi cover and the puncture-opening
map of \cite[Section 2]{CA_NT2002}.  Directly,
\[
 1-q(t)=\frac{(t+1)^2(2t-1)^2(t-2)^2}
                    {4(t^2-t+1)^3}.
 \tag{B.C8}
\]
The six anharmonic transformations preserve $q$, whose
degree is six.  Its local degrees over $0,1,\infty$ are,
respectively, two, two, and three.  These computations prove
that $q$ is the asserted unramified $\mathfrak S_3$ orbifold
cover.

Choose a small rational $t_0>0$, and base the target at
$T_0=q(t_0)$.  Take positive meridians $\tau$ around $T=0$
and $s$ around $T=1$, with stems along the real interval.
Order the peripheral product as
$\ell_1\ell_0\ell_\infty=1$ and set $r=s\tau$.
This gives $s^2=r^3=1$ and $\tau=sr$ in the orbifold group.
The positive meridian $y_L$ about $t=0$ has image $\tau^2$:
the local expansion is $q(t)=(27/4)t^2+O(t^3)$.
The lift of $s$ from $t_0$ ends at $1-t_0$, passing around
the branch point $1/2$ and implementing $t\mapsto1-t$.
Transport the positive meridian about $t=1$ back along this
lift to define $x_L$.  The holomorphic deck involution
preserves orientation.  We obtain the actual inclusion
\[
             q_*(x_L)=(rs)^2,\qquad q_*(y_L)=(sr)^2.
 \tag{B.C9}
\]
Under $j$, the detour around $1/2$ contracts because that is
an ordinary target point.  Transport the basepoint from
$T_0$ to $t_0$ along the real interval.  These same source
loops then have $j_*(x_L)=s$ and $j_*(y_L)=\tau$.
Thus (B.C9) and the physical quotient in (B.C2) occur in
one simultaneous oriented marking.  Both algebraic maps in
(B.C7) are defined over $\mathbb Q$, so their arithmetic
fundamental-group extensions commute with the Galois outer
actions.  A basepoint path changes this statement only by a
geometric inner automorphism, already removed in descent and
restriction.  This proves that $D_L$ fixes the arithmetic
image in the LM2 basis.

Changing that basis by $\kappa=\sigma c$ preserves the
arithmetic subgroup: $\sigma$ commutes with its outer action,
and $c$ belongs to it.  Hence the pointwise fixed assertion
holds for every element of the same arithmetic subgroup in
the current basis as well.  Equations (B.C3) now give (B.C6).
To compute its centralizer of $c$, observe that an
automorphism commuting with $c$ preserves
$\mathbb R_{\rm alg}=\overline{\mathbb Q}\cap\mathbb R$.
The positive elements there are exactly the nonzero squares.
It therefore preserves the order, fixes the dense subfield
$\mathbb Q$, and fixes each real algebraic number.  As
$\overline{\mathbb Q}=\mathbb R_{\rm alg}(i)$, it is $1$
or $c$.  This proves the exact intersection in (B.C6).

For the explicit finite witness put
$K=\mathbb Q(\zeta_8)$, $a=2^{1/8}>0$, and $L=K(a)$.
The real quartic field generated by $a^2$ is not normal:
$T^4-2$ is Eisenstein, and that field misses the nonreal root
$ia^2$.  Thus it is not a subfield of the abelian extension
$K/\mathbb Q$.  Since $\mu_4\subset K$ and $a^4=\sqrt2$
belongs to $K$, the Galois group of $K(a)/K$ is a subgroup
of $\mu_4$.  Degree at most two would make $a^2\in K$.
Consequently $[L:K]=4$, $[L:\mathbb Q]=16$, and $L$ is the
splitting field of $T^8-2$.

Write an automorphism as $(v,b)$, for
$\zeta_8\mapsto\zeta_8^v$, $a\mapsto\zeta_8^ba$.
Its exact set and multiplication are
\[
\begin{gathered}
 v\in\{1,3,5,7\},\quad b\in\mathbb Z/8,\qquad
 b\equiv0\pmod2\ (v=1,7),\quad
 b\equiv1\pmod2\ (v=3,5),\\
       (v,b)(v',b')=(vv',b+vb').
\end{gathered}                                             \tag{B.C10}
\]
The parity condition follows from
$a^4=\zeta_8+\zeta_8^{-1}$; there are exactly sixteen such
pairs, so they give the full group.  Its derived group is
$\{(1,0),(1,2),(1,4),(1,6)\}$, and explicitly
\[
 [(3,1),(5,3)]=(1,2),\qquad
 (7,0)(1,2)(7,0)=(1,6).
 \tag{B.C11}
\]
Lift the first two factors to $G_{\mathbb Q}$ and take their
commutator $\gamma$.  It lies in
$G_{\mathbb Q^{\rm ab}}$ and has image $(1,2)$.
Complex conjugation has image $(7,0)$, so (B.C11) and (B.C6)
distinguish $D_P(\gamma)$ from $\gamma$.  This is an actual
arithmetic element, not a finite pair for which a GT lift is
being assumed.  R1 gives
$\gamma\in\widehat{GT}$ and $\Delta(\gamma)=0$.
It proves all three refutations stated in the result, while
also proving $C_L\ne C_P$.
\end{proof}

\paragraph{Exact multiplier-one descent formulas.}
The following formulas will also identify the kernel element.
For $(1,f)\in G$ put
\[
 C=f(s,\tau),\quad H=f(r^{-1},s),\quad
 C^{\rm ab}=(p,q),\quad a_0=q-p,\quad H\bmod F=r^l.
 \tag{B.C12}
\]
Then $C\in F$, and with
$\sigma=\operatorname{Ad}_s|_F$ the normalized pair and pure
frames of $D_P(1,f)$ are
\[
 f_{D_P}=y^{a_0}\sigma(C)C^{-1}x^{-a_0},\quad
 \beta_{D_P}=x^{a_0}C,\quad
 A_{D_P}=x^{a_0}CHr^{-l},\quad \rho_{D_P}=-q.
 \tag{B.C13}
\]
To check these statements before normalizing, the actual
descendant is
$s\mapsto s$, $\tau\mapsto C^{-1}\tau C$,
$r\mapsto HrH^{-1}$.  The last equality follows from the
third-peripheral formula and $\varphi_P(z)=r^{-1}$.
Both $C,H$ map into $A_3$ in $\mathfrak S_3$ because $f$ is
a commutator word.  The descendant fixes $s,r$ there; hence
$C\bmod F$ centralizes the transposition $\tau$.  Its
intersection with $A_3$ is trivial, proving $C\in F$.
Conjugating by $C$ fixes $\tau$ and gives pair coordinate
$\sigma(C)C^{-1}$ of abelianization $(a_0,-a_0)$.
The further conjugation by $x^{a_0}$ gives exactly (B.C13).
Its harmonic abelianization is $(q,q)$, proving the last
formula.  The pure cubic frame is the displayed one because
$r^l$ centralizes $r$; its zero abelianization also follows
from $(1-\operatorname{Ad}_r)A^{\rm ab}=0$ and the
injectivity of this matrix, of determinant three, over
$\widehat{\mathbb Z}$.

For comparison, let $a=\tau^{-1}$,
$P=f(a,s)\in F$, $P^{\rm ab}=(p_L,q_L)$, and
$b_0=p_L-q_L$.  The same calculation for $\varphi_L$ gives
\[
 \beta_{D_L}=x^{b_0}P^{-1},\qquad
 (1,f)\in C_L\ \Longleftrightarrow\
          \beta^{-1}=f(a,s)x^{-b_0}.
 \tag{B.C14}
\]
Indeed the initial descendant fixes $a$, sends
$s$ to $P^{-1}sP$, and is the identity modulo the pure
subgroup.  The unique conjugation by an $x$-power making
its pair a commutator is $x^{b_0}$.  A normalized action
fixing $\tau$ is determined by its pure $s$-frame, proving
the equivalence.  Equations (B.C13)--(B.C14) retain the exact
morphism between the two primitive formulas.

\begin{result}{B.1.3}{Proved}{A full profinite kernel element detected by the CM character}
There exists an actual multiplier-one element $k\in G$ with
\[
                   D_P(k)=D_L(k)=1,\qquad\Delta(k)=1.
 \tag{B.C15}
\]
It can be chosen with
\[
 f_k,\beta_k\in R_{\rm nil}(U_0),\quad
 U_0=F'F^3,\quad \rho_k=0,\quad
 q_{12}(\beta_k)=1,\quad q_{12}(A_k)=[5I].
 \tag{B.C16}
\]
Here $R_{\rm nil}(U)$ is the kernel of the maximal
pronilpotent quotient of the actual profinite group $U$,
and $q_{12}:B\twoheadrightarrow
\operatorname{SL}_2(\mathbb Z/12)/\{\pm I\}$ is the fixed
modular quotient.  This proves, in particular, that neither
$C_P$ nor $C_L$ contains $\ker\Delta$.
\end{result}
\begin{proof}
We verify the actual source of the relative lifting theorem,
so the argument does not infer a lift from finite survival.
The word
\[
                       A_0=[[x^4,y^4],z^4]^2              \tag{B.C17}
\]
belongs to $F'\cap\operatorname{core}_B(\ker\varphi_P)$.
There are six kernels to check because $F$ is normal of
index six and $\ker\varphi_P$ is normal in $F$.
The six deck substitutions permute the three peripheral
conjugacy classes.  Under each followed by $\varphi_P$,
one of $x,y,z$ becomes a conjugate of $s$, whose fourth
power is one.  An iterated commutator in (B.C17) is one
whenever any of its three inputs is one.  This proves
membership in all six full profinite kernels.

For its original modular value use the actual matrices
$x\mapsto U^2$, $y\mapsto L^2$, where
$U=\left(\begin{smallmatrix}1&1\\0&1\end{smallmatrix}\right)$,
$L=\left(\begin{smallmatrix}1&0\\-1&1\end{smallmatrix}\right)$.
Multiplication modulo twelve gives $A_0\mapsto5I$.
The image $q_{12}(U_0)$ has order $32$ and is a two-group;
it is also described in R1 as $V_4\times Q_8$.
The independent exact program
\texttt{check\_descent\_audit.py} enumerates this image
from commutators and cubes of the actual matrix image of
$F$, and checks $[5I]$ belongs to it.  These finite
computations are used only for this source quotient.

For each open normal $J\triangleleft B$ put
\[
 N_J=U_0\cap\operatorname{core}_B(\varphi_P^{-1}(J)),\quad
 D_J=B/N_J,\quad
 E_J=\operatorname{im}\bigl(B\longrightarrow Q_{12}\times D_J\bigr).
 \tag{B.C18}
\]
These are actual finite quotients; $N_J$ is open normal
in $B$, is contained in $F$, and contains $A_0$.
The kernel of $E_J\to D_J$ is the two-group $q_{12}(N_J)$.
Consequently $q_{12}$ factors through
$\Gamma_2=B/O^2(N_J)$, where $O^2(N_J)$ is the kernel of
the maximal pro-two quotient of $N_J$, and
$\Gamma_2\twoheadrightarrow E_J$ has pro-two kernel.
The element $([5I],1)$ belongs to $E_J$, by (B.C17), and
centralizes the image of $r$, of order three.  Coprime
centralizer lifting gives
\[
 c_2\in C_{\Gamma_2}(r),\qquad c_2\longmapsto([5I],1).
 \tag{B.C19}
\]
Explicitly, for any initial lift of this element, its
conjugate of $\langle r\rangle$ and $\langle r\rangle$
are two complements to the same pro-two kernel over the
same order-three subgroup.  Schur--Zassenhaus at finite
quotients conjugates them by the kernel.  The generator
has the same order-three quotient image, so the adjusted
lift centralizes $r$ itself.  The closed nonempty sets of
finite adjustments have the finite intersection property;
compactness supplies a single profinite adjustment.

Use the actual relative product in the refereed R1/D.50
framed lifting theorem,
\[
 B/R_{\rm nil}(N_J)\simeq
              \prod_{p,D_J} B/O^p(N_J).                  \tag{B.C20}
\]
Place the identity automorphism in every component.
Choose the positive $s$ and cusp conjugators to be one
in every component.  Choose the positive $r$ conjugator
to be $c_2$ in the two-component and one elsewhere.
They are valid conjugators for these identity actions by
(B.C19).  Their images in $D_J$ are all one.  These facts
verify the exact common-quotient and peripheral hypotheses
of the framed theorem, which gives an actual full
peripheral automorphism with the prescribed coordinates.
Its lifted cusp conjugator lies in $R_{\rm nil}(N_J)$;
conjugating by its inverse normalizes the cusp without
changing any retained finite values.  The resulting
normalized gentle element $k_J$ has
\[
 f_J,\beta_J\in R_{\rm nil}(N_J),\quad
 q_{12}(\beta_J)=1,\quad q_{12}(A_J)=[5I].
 \tag{B.C21}
\]
Since $R_{\rm nil}(N_J)\subseteq N_J'\subseteq F'$, its
harmonic frame has zero abelianization, so $\rho_J=0$.
The normalized finite-frame formula of R1 gives
$\Delta(k_J)=1$.  Moreover $\varphi_P(f_J)\in J$.

Pass to the actual compact multiplier-one group, imposing
the closed conditions
\[
\begin{gathered}
 f,\beta\in R_{\rm nil}(U_0),\quad
 q_{12}(\beta)=1,\quad q_{12}(A)=[5I],\\
 \varphi_P(bfb^{-1}),\varphi_P(b\beta b^{-1}),
 \varphi_P(bAb^{-1})\in J,\\
 b\in\{1,s,r,r^2,sr,sr^2\}.
\end{gathered}                                         \tag{B.C22}
\]
The normalized pair and its unique pure frames are
continuous in this group.  Every finite collection of
these conditions is met by the construction with the
intersection of its $J$'s: for $N_J\subseteq U_0$,
$R_{\rm nil}(N_J)\subseteq R_{\rm nil}(U_0)$, since
every finite nilpotent quotient of $U_0$ restricts to
one of $N_J$.  Each $f_J,\beta_J$ also belongs to $N_J$.
The retained cubic conjugator has common $D_J$-value one,
so $A_J\in N_J$ as well.  Thus all three satisfy the six
displayed core conditions.
Compactness therefore produces a single actual $k$ with
(B.C16) and $f_k,\beta_k,A_k\in\operatorname{core}_B\ker\varphi_P$.
In (B.C13) this means $C=1$, $p=q=a_0=0$, and
$\beta_{D_P}=1$.  The normalized descendant fixes both
$s$ and $\tau$, so it is the identity.  Equation (B.C3)
also gives $D_L(k)=1$.  The retained frames give
$\Delta(k)=1$, proving (B.C15).

The element $k$ is nonidentity.  R1's multiplier-one
torsion-freeness gives $k^2\ne1$, whereas
$\Delta(k^2)=0$ and $D_P(k^2)=D_L(k^2)=1$.
Thus $k^2$ belongs to $\ker\Delta$ and belongs to neither
fixed subgroup, proving the final assertion and refuting
the proposed equality $C_L=\ker\Delta$ as well.
\end{proof}

\paragraph{What is decided about further quadratic characters.}
For $n\geq1$ the actual continuous characters
$\chi_n=\Delta\circ D_P^n$ vanish on $G_{\mathbb Q}$
by (B.C6).  They all vanish at $k$, whereas $\Delta(k)=1$.
Thus $\Delta$ is different from each $\chi_n$ and is not
in their finite $\mathbb F_2$-linear span.  More precisely,
\[
 (D_P,\Delta):G\twoheadrightarrow
                    \operatorname{im}(D_P)\times\mathbb F_2
 \tag{B.C23}
\]
is surjective: $g$ and $gk$ have the same descendant and
opposite bits.  This also holds on the multiplier-one
group.  B.2.8 now proves $\chi_1\ne0$ by an actual full
finite-certificate lift, and hence proves that $\Delta$
and $\chi_1$ are independent. B.2.24--B.2.25 subsequently
prove surjectivity of $D_P$ and independence of every
positive iterate, by actual triangular witnesses.

There is a proved partial restriction on both fixed
subgroups.  In the actual affine quotient (A.J18), the
translation of $P=f(\tau^{-1},s)$ is $ct_f$, while the
translation of $C=f(s,\tau)$ is $-ct_f$; here
$t_f=\overline{t_f}$ by (A.J17).  Taking inverses of
the two frames in (B.C13)--(B.C14) and multiplying by
their zero-translation correcting shears gives
\[
 \epsilon_2(W_{D_P g})=\epsilon_2(W_{D_L g})
                  =\Psi_3(g)/2\pmod2.                  \tag{B.C24}
\]
All rightmost $x$-powers have zero translation, so no
uncomputed normalization exponent enters this equality.
Consequently $d_2(g)=0$ on $C_P\cap\mathcal H$ and
$C_L\cap\mathcal H$.  On the full gentle multiplier-one
domain define
$\widetilde e_3=\epsilon_3(T)+\Psi_3/2\pmod2$; it agrees
with $e_3$ on the five-point domain where $\kappa_d$ is
defined.  The exact identity
$\Delta=d_2+\widetilde e_3$ makes the remaining comparison
on either fixed subgroup $\widetilde e_3=0$.
Neither the counterexamples above
nor (B.C24) proves or refutes $C_P\subseteq\ker\Delta$
or $C_L\subseteq\ker\Delta$.  The physical inclusion $\widehat{GT}\subseteq C_L$ is now proved
by the cyclic quotient and equianharmonic opening comparison
in (B.K1)--(B.K13).  Thus every positive iterate
$\Delta\circ D_P^n$ vanishes on full GT.  Their nonvanishing
and the remaining fixed-subgroup kernel inclusions are not
inferred from the finite source census.
